\section{Conclusion}
The relevant possibility is a growing civilization, not necessarily a
satiated or inactive one. Internal computational uses, including recursively
generated virtual worlds, can compete successfully with additional external
projects. The resource frontier and the common-budget model make the two
requirements distinct: the chosen workload must remain physically
supportable, and inward investment must remain preferable at the relevant
opportunity costs.

Continued internal growth with no additional external investment is
therefore conditional on a path, not established by a favourable corner
at one budget. Under logarithmic computational value, budgets tending to
infinity, fixed positive residual external slopes and bounded relative
prices are incompatible with persistent all-inward allocation, irrespective
of productivity. Persistence requires the price-to-residual-slope ratio to
exceed $(1+B)\ln(1+B)$ wherever a residual slope is positive, together
with the exact productivity threshold and continuing local feasibility.
The curvature extension further shows that the finite residual floor is
specific to logarithmic value: slower marginal-value decay permits an
eventual productivity threshold at fixed budgets and prices, whereas
faster decay makes sufficiently high productivity incompatible with
inwardness when residual slopes are positive.

Concave external alternatives preserve the corner comparison when initial
returns are unchanged, but alter mixed allocations. Cheap external access
and future resource production supply additional reasons to expand.
These boundaries define a quantitatively specified alternative to
inferring inevitable expansion from capability alone.
They do not establish its prevalence, indefinite engineering realization,
or a population-level explanation of galactic silence.

